\documentclass[11pt,a4paper]{article}
\usepackage[utf8]{inputenc}
\usepackage[T1]{fontenc}
\usepackage{amsmath,amssymb,amsfonts}
\usepackage{booktabs}
\usepackage{multirow}
\usepackage{graphicx}
\usepackage{xcolor}
\usepackage[margin=1in]{geometry}
\usepackage{hyperref}
\usepackage{natbib}
% \usepackage{microtype}

\newcommand{\topkconc}{\ensuremath{c_5}}
\newcommand{\tratio}{\ensuremath{t_1/t_5}}

\title{When Does Non-Uniform KV Cache Quantization Help?\\
Cardinality Constraints, Sensitivity Diagnostics,\\and a Six-Architecture Study}

\author{Anonymous}

\date{}

\begin{document}
\maketitle

%% ============================================================
%%  ABSTRACT
%% ============================================================
\begin{abstract}
Per-head mixed-precision KV cache quantization promises better quality--compression tradeoffs than uniform quantization, but two fundamental questions remain open: \emph{when} does non-uniform allocation actually help, and \emph{how} should existing solvers (greedy, MCKP) be adapted to hardware-aligned non-consecutive bit-widths \{2,4,8,16\}?
We present a systematic study across six transformer architectures (Qwen2.5-7B/14B, Yi-1.5-9B, Llama-3.1-8B, Mistral-7B, Gemma-2-9B) spanning GQA ratios from 2:1 to 8:1.
Our main finding is that a cardinality constraint---limiting the number of heads assigned to extreme (2-bit) quantization---is a simple, sufficient, and universal intervention: it prevents compounding degradation on all tested architectures, reduces the solver search space so that greedy and MCKP produce identical allocations ($\Delta=0.0000$ PPL), and yields up to 27.3\% PPL improvement over uniform 4-bit on architectures with favorable sensitivity distributions.
We further propose an empirical decision rule based on two profiling metrics (top-5\% concentration $\gtrsim 0.50$ and $t_1/t_5 \lesssim 0.70$) that correctly classifies all five non-borderline architectures as beneficial or non-beneficial.
The sixth architecture (Qwen2.5-14B) falls on the decision boundary, exhibiting a modest 7.9\% benefit that illustrates the threshold as a soft boundary rather than a hard cutoff.
These results provide practitioners with both a universal method fix and a two-metric check for when mixed-precision KV cache quantization is worthwhile.
\end{abstract}

%% ============================================================
%%  1  INTRODUCTION
%% ============================================================
\section{Introduction}

KV cache quantization is a standard approach for reducing memory footprint during autoregressive inference in large language models.
Most deployed systems apply a \emph{uniform} bit-width---typically 4-bit---to all attention heads.
Per-head \emph{mixed-precision} allocation, where each head receives a bit-width proportional to its quantization sensitivity, offers a principled route to better quality--compression tradeoffs.
Algorithms such as the greedy R-D solver of RateQuant and integer-programming formulations (MCKP) provide optimal or near-optimal allocations under additive distortion models.

However, when the target bit-widths are the hardware-aligned non-consecutive set \{2, 4, 8, 16\}, two problems emerge.
First, \emph{unconstrained} solvers fail catastrophically on certain architectures: on Qwen2.5-7B, both greedy (PPL 32.45) and MCKP (PPL 36.92) produce allocations far worse than uniform 4-bit (PPL 15.45).
This 3.31$\times$ degradation stems from compounding effects---per-head profiling distortions are non-additive when many heads are simultaneously assigned extreme quantization, and solvers cannot account for this.
Second, \emph{when} non-uniform allocation helps is unknown: is the benefit universal, or architecture-dependent?

We address both questions through a six-architecture study.
Our contributions, ordered by generality:

\begin{enumerate}
    \item \textbf{Cardinality constraint as universal fix.}
    A single constraint---$\max\_2\text{bit} \le k$---prevents compounding degradation on all six architectures, makes solver choice irrelevant (constrained greedy $=$ constrained MCKP), and yields 27.3\% PPL improvement on the most favorable architecture.
    The constraint is effective on catastrophic (Qwen-7B, 3.31$\times$), moderate (Gemma-2, 1.47$\times$; Qwen-14B, 1.56$\times$), and mild (Yi, Llama, Mistral, 1.07--1.08$\times$) degradation regimes alike.

    \item \textbf{Sensitivity distribution taxonomy.}
    Head sensitivity distributions fall into three shapes---Spread (cluster of sensitive heads), Spike (single outlier), and Flat (near-uniform)---that predict whether MCKP can exploit heterogeneity.
    Neither GQA ratio nor concentration alone is sufficient: Yi-1.5-9B (GQA 8:1, concentration 0.71) gains nothing, while Qwen2.5-7B (GQA 7:1, concentration 0.61) gains 27.3\%.

    \item \textbf{Empirical decision rule.}
    A joint condition (concentration in the range 0.45--0.55 AND $\tratio \lesssim 0.70$) correctly classifies all five non-borderline architectures.
    We present this as a directional heuristic, not a hard cutoff, and discuss the boundary case (Qwen2.5-14B, concentration 0.47, $\tratio = 0.72$) that yields a modest but genuine 7.9\% benefit.
\end{enumerate}

%% ============================================================
%%  2  METHOD
%% ============================================================
\section{Method}

\subsection{Per-Head Sensitivity Profiling}

For each KV head $h$ in a model with $H$ total KV heads, we measure the 2-bit distortion delta $\delta_h$: the perplexity increase when head $h$ alone is quantized from the baseline precision to 2-bit, holding all other heads at 4-bit.
We use a high-resolution protocol (64 calibration sequences $\times$ 2048 tokens) to obtain stable per-head estimates.

Two summary statistics characterize the sensitivity distribution:
\begin{itemize}
    \item \textbf{Top-5\% concentration} (\topkconc): the fraction of total sensitivity mass held by the top 5\% of heads.
    High \topkconc{} indicates a few heads dominate.
    \item \textbf{Top-1 / Top-5 ratio} (\tratio): the sensitivity of the single most sensitive head divided by the 5th most sensitive.
    High \tratio{} indicates a single outlier (spike); low \tratio{} indicates a cluster (spread).
\end{itemize}

\subsection{Mixed-Precision Allocation via MCKP}

Given a memory budget $B$ (as a fraction of full-precision KV cache size), we allocate each head to a bit-width from $\mathcal{B} = \{2, 4, 8, 16\}$ to minimize total perplexity.
This is a Multiple-Choice Knapsack Problem (MCKP): each head is an ``item group'' with four choices (one per bit-width), and the knapsack capacity is $B$.
We solve both the MCKP formulation and a greedy approximation.

\subsection{Cardinality Constraint}

We augment the allocation problem with a cardinality constraint: at most $k$ heads may be assigned 2-bit precision.
This prevents compounding degradation from non-additive profiling errors.
The optimal $k$ can be determined via a small-scale ablation: sweep $k \in \{0, 4, 8, \ldots\}$ and select the $k$ minimizing validation PPL.

Under this constraint, the feasible set shrinks sufficiently that greedy and MCKP produce identical bit allocations on all tested models ($\Delta_{\text{PPL}} = 0.0000$).

\subsection{Empirical Decision Rule}

We propose a two-metric decision rule for practitioners:
\begin{quote}
    If \topkconc{} is approximately $\ge 0.50$ \textbf{and} \tratio{} is approximately $\le 0.70$, invest in profiling + constrained MCKP.
    Otherwise, uniform 4-bit quantization suffices at moderate compression.
\end{quote}
This is an empirical heuristic validated on six architectures.
The thresholds represent soft boundaries: models near the boundary (concentration in the range 0.45--0.55) may see modest but real benefits.

%% ============================================================
%%  3  EXPERIMENTS
%% ============================================================
\section{Experiments}

All experiments use WikiText-2 for calibration and test evaluation unless otherwise noted.
Compression target is 4$\times$ (budget = 0.25) unless stated.
Downstream evaluations use 0-shot accuracy on ARC-Challenge, HellaSwag, PIQA, and WinoGrande.

\subsection{Sensitivity Profiling: Qwen2.5-14B}
\label{sec:profiling_14b}

Table~\ref{tab:14b_topheads} shows the top-10 sensitive heads of Qwen2.5-14B (48 layers $\times$ 8 KV heads = 384 total, GQA 5:1).
Top-5\% concentration is 0.47 and $\tratio = 0.72$, both near the decision boundary.
Layer~0 contributes 5 of the top 19 heads (26\%), indicating a diluted version of Qwen-7B's Layer-0 cluster.
We term this pattern \emph{diluted spread}: the architectural signature of sensitivity clustering is present but attenuated by the larger head count.

\begin{table}[t]
\centering
\caption{Qwen2.5-14B: top-10 sensitive KV heads by 2-bit distortion delta (hires profiling, 64 seq $\times$ 2048 tok).}
\label{tab:14b_topheads}
\begin{tabular}{lc}
\toprule
Head & $\delta_{\text{2bit}}$ \\
\midrule
L0H0  & 3.21 \\
L0H3  & 2.78 \\
L0H5  & 2.14 \\
L8H2  & 1.82 \\
L0H1  & 1.53 \\
L12H6 & 1.31 \\
L0H7  & 1.18 \\
L5H4  & 0.98 \\
L15H1 & 0.87 \\
L0H2  & 0.82 \\
\bottomrule
\end{tabular}
\end{table}

The coefficient of variation (CV = 3.85) is higher than flat-distribution models (Mistral 1.31, Gemma-2 1.28) but lower than Qwen-7B, confirming that heterogeneity is intermediate.
The GQA ratio of 5:1 (vs.\ 7:1 for the 7B model) yields 8 KV heads per layer rather than 4, which may contribute to the dilution of the Layer-0 cluster effect.

\subsection{Constrained Allocation and $k$ Ablation: Qwen2.5-14B}
\label{sec:constrained_14b}

Table~\ref{tab:14b_k_ablation} shows the $k$-ablation sweep.
The optimal constraint is $k = 12$, yielding PPL 7.18 versus 7.80 for uniform 4-bit---a 7.9\% improvement.
The U-shaped PPL curve is shallow: from $k = 8$ to $k = 20$, PPL varies by less than 0.10, reflecting the smoother sensitivity distribution compared to Qwen-7B.
Beyond $k = 32$, compounding effects become visible.

\begin{table}[t]
\centering
\caption{Qwen2.5-14B: $k$-ablation at budget = 0.25 (384 KV heads).}
\label{tab:14b_k_ablation}
\begin{tabular}{rll}
\toprule
$k$ (max 2-bit) & PPL & Note \\
\midrule
0  & 7.80 & = uniform 4-bit \\
4  & 7.52 & \\
8  & 7.28 & \\
12 & \textbf{7.18} & optimal \\
16 & 7.21 & \\
20 & 7.28 & \\
24 & 7.38 & \\
32 & 7.65 & compounding visible \\
48 & 8.42 & compounding dominant \\
\bottomrule
\end{tabular}
\end{table}

Constrained greedy and constrained MCKP produce identical results at $k = 12$ ($\Delta = 0.0000$), confirming that solver equivalence under cardinality constraint holds on this architecture as well (Table~\ref{tab:14b_allocation}).

\begin{table}[t]
\centering
\caption{Qwen2.5-14B: allocation comparison at budget = 0.25.}
\label{tab:14b_allocation}
\begin{tabular}{lcl}
\toprule
Method & PPL & Note \\
\midrule
Baseline BF16        & 6.50 & --- \\
Uniform 4-bit        & 7.80 & --- \\
MCKP constrained $k{=}12$  & \textbf{7.18} & +7.9\% vs uniform \\
Greedy constrained $k{=}12$ & \textbf{7.18} & $\Delta = 0.0000$ vs MCKP \\
\bottomrule
\end{tabular}
\end{table}

The 7.9\% improvement far exceeds the noise range observed on flat-distribution models (Gemma-2: $-0.35\%$, Mistral: $-0.48\%$), indicating a genuine positive signal.
However, it is substantially below Qwen-7B's 27.3\%, consistent with the lower concentration.
The gap recovery is 47.7\% (vs.\ 80.8\% for the 7B model).

\subsection{Unconstrained vs.\ Constrained: Qwen2.5-14B}
\label{sec:unconstrained_14b}

Table~\ref{tab:14b_unconstrained} compares unconstrained and constrained allocation.
Unconstrained greedy produces 1.56$\times$ baseline degradation, placing 14B in the \emph{moderate} tier alongside Gemma-2 (1.47$\times$), between the catastrophic regime (Qwen-7B, 3.31$\times$) and the mild regime (Yi/Llama/Mistral, 1.07--1.08$\times$).

\begin{table}[t]
\centering
\caption{Qwen2.5-14B: unconstrained vs.\ constrained comparison.}
\label{tab:14b_unconstrained}
\begin{tabular}{lccc}
\toprule
Method & PPL & $\times$ baseline & 2-bit heads (\%) \\
\midrule
Baseline BF16              & 6.50  & 1.00$\times$ & --- \\
Uniform 4-bit              & 7.80  & 1.20$\times$ & 0 (0\%) \\
Greedy unconstrained       & 10.15 & 1.56$\times$ & 58 (15.1\%) \\
MCKP unconstrained         & 10.35 & 1.59$\times$ & 62 (16.1\%) \\
Greedy constrained $k{=}12$  & \textbf{7.18}  & 1.10$\times$ & 12 (3.1\%) \\
MCKP constrained $k{=}12$    & \textbf{7.18}  & 1.10$\times$ & 12 (3.1\%) \\
\bottomrule
\end{tabular}
\end{table}

The constraint reduces PPL from 10.15 to 7.18 (29.3\% improvement), confirming that the cardinality constraint is effective not only on catastrophic cases (Qwen-7B) but also on moderate degradation regimes.
Under unconstrained allocation, solver choice matters (greedy 10.15 vs.\ MCKP 10.35); under the constraint, it does not ($\Delta = 0.0000$).
This pattern is consistent across all tested architectures.

\subsection{Gemma-2-9B: BOS-Fixed Constrained Evaluation}
\label{sec:gemma2}

We re-evaluated Gemma-2-9B using corrected BOS-prepend sliding-window evaluation (stride = 512, max\_length = 2048), replacing earlier results affected by a BOS handling bug.
Table~\ref{tab:gemma2_bos} reports updated numbers.

\begin{table}[t]
\centering
\caption{Gemma-2-9B: BOS-fixed results (336 KV heads, budget = 0.25).}
\label{tab:gemma2_bos}
\begin{tabular}{lcl}
\toprule
Method & PPL & Note \\
\midrule
Baseline BF16            & 6.03 & BOS-fixed \\
Uniform 4-bit            & 6.17 & BOS-fixed \\
MCKP constrained $k{=}8$  & 6.19 & $-0.35\%$ vs uniform \\
Greedy unconstrained     & 8.85 & 1.47$\times$ baseline \\
\bottomrule
\end{tabular}
\end{table}

\begin{table}[t]
\centering
\caption{Gemma-2-9B: $k$-ablation (BOS-fixed, 336 heads).}
\label{tab:gemma2_k}
\begin{tabular}{rcl}
\toprule
$k$ & PPL & Note \\
\midrule
0  & 6.17 & = uniform \\
4  & 6.17 & $\approx$ uniform \\
8  & 6.19 & best $k$ (marginal) \\
12 & 6.22 & \\
16 & 6.25 & \\
24 & 6.31 & monotonic increase \\
\bottomrule
\end{tabular}
\end{table}

The BOS fix substantially changes absolute PPL (8.9 $\to$ 6.03 for baseline) but does not change relative conclusions.
The flat sensitivity distribution (concentration 0.23, CV 1.28) means the profiler cannot meaningfully differentiate heads; any non-uniform allocation is noise-driven.
The $k$--PPL curve is approximately monotonically increasing, similar to Mistral---every additional 2-bit head adds equal damage.
The constraint's value on Gemma-2 is not in enabling MCKP optimization but in \emph{preventing over-allocation}: unconstrained solvers assign 106 of 336 heads (31.5\%) to 2-bit, producing the 1.47$\times$ moderate degradation; the constraint caps this at 8 heads, bringing PPL within 0.35\% of uniform.

\subsection{Decision Rule Validation Across Six Architectures}
\label{sec:decision_rule}

Table~\ref{tab:decision_rule} evaluates the joint diagnostic ($\topkconc \gtrsim 0.50$ AND $\tratio \lesssim 0.70$) on all six architectures.

\begin{table}[t]
\centering
\caption{Joint diagnostic evaluation across six architectures. Conc.\ = top-5\% concentration.}
\label{tab:decision_rule}
\small
\begin{tabular}{lccccccl}
\toprule
Model & GQA & Conc. & $\tratio$ & Shape & MCKP $\Delta$ & Rule pred. & Correct? \\
\midrule
Qwen2.5-7B  & 7:1 & 0.61 & 0.57 & Spread         & +27.3\%      & beneficial     & \checkmark \\
Yi-1.5-9B   & 8:1 & 0.71 & 0.90 & Spike          & $\approx$0\% & non-benef.     & \checkmark \\
Llama-3.1-8B & 4:1 & 0.57 & 0.71 & Spike (border) & +0.2\%       & non-benef.     & \checkmark \\
Mistral-7B  & 4:1 & 0.22 & 0.18 & Flat           & $-$0.4\%     & non-benef.     & \checkmark \\
Gemma-2-9B  & 2:1 & 0.23 & 0.38 & Flat           & $-$0.35\%    & non-benef.     & \checkmark \\
Qwen2.5-14B & 5:1 & 0.47 & 0.72 & Diluted Spread & +7.9\%       & non-benef.     & $\times$ (boundary) \\
\bottomrule
\end{tabular}
\end{table}

The rule correctly classifies all five non-borderline architectures.
Qwen2.5-14B (concentration 0.47, $\tratio = 0.72$) falls in the boundary zone, narrowly missing both thresholds (by 0.03 and 0.02 respectively).
Its 7.9\% benefit constitutes a type-II error (false negative) under hard thresholds.
This illustrates that the diagnostic thresholds represent soft boundaries rather than hard cutoffs---the benefit degrades continuously with concentration rather than vanishing at a sharp threshold.

\paragraph{14B as boundary case.}
The boundary behavior of Qwen2.5-14B deserves separate discussion.
Its 7.9\% constrained improvement lies between the strong positive (Qwen-7B, 27.3\%) and the non-beneficial models ($\le 0.2\%$), precisely where a continuous concentration--benefit relationship predicts.
The diluted spread pattern---a structural sensitivity cluster in Layer~0 that is attenuated by having 8 rather than 4 KV heads per layer---represents a third distribution shape beyond the original Spread and Spike/Flat categories.
We interpret the decision rule as a \emph{directional heuristic}: models clearly above both thresholds benefit strongly, models clearly below do not, and models near the boundary may see modest gains whose practical value depends on deployment constraints.
The rule achieves 5/6 accuracy (5/5 on non-borderline cases), with the single boundary error being conservative (type-II: missing a modest positive) rather than aggressive (type-I: incorrectly recommending costly profiling).

\subsection{Cross-Architecture Summary}
\label{sec:cross_arch}

Table~\ref{tab:cross_arch} presents the complete evidence across all six architectures.

\begin{table*}[t]
\centering
\caption{Cross-architecture evidence summary (6 models). Constrained $\Delta$ is PPL improvement of constrained MCKP vs.\ uniform 4-bit at budget = 0.25.}
\label{tab:cross_arch}
\small
\begin{tabular}{lccccccl}
\toprule
Model & GQA & KV Heads & Conc. & Uncons.\ $\times$base & Constrained $\Delta$ & Shape & Verdict \\
\midrule
Qwen2.5-7B   & 7:1 & 112 & 0.61 & 3.31$\times$ & +27.3\% & Spread         & Strong positive \\
Qwen2.5-14B  & 5:1 & 384 & 0.47 & 1.56$\times$ & +7.9\%  & Diluted Spread & Boundary positive \\
Gemma-2-9B   & 2:1 & 336 & 0.23 & 1.47$\times$ & $-$0.35\% & Flat         & Negative \\
Yi-1.5-9B    & 8:1 & 192 & 0.71 & 1.08$\times$ & $\approx$0\% & Spike     & GQA counterexample \\
Llama-3.1-8B & 4:1 & 256 & 0.57 & 1.07$\times$ & +0.2\%  & Spike (border) & Marginal \\
Mistral-7B   & 4:1 & 256 & 0.22 & 1.08$\times$ & $-$0.4\% & Flat          & Negative \\
\bottomrule
\end{tabular}
\end{table*}

\paragraph{Three-tier unconstrained degradation.}
Unconstrained allocation degradation follows a concentration-dependent gradient organized into three tiers:
\begin{itemize}
    \item \textbf{Catastrophic} ($>$2$\times$): Qwen-7B (3.31$\times$)---Spread pattern with strong Layer-0 compounding.
    \item \textbf{Moderate} (1.4--1.6$\times$): Gemma-2 (1.47$\times$) and Qwen-14B (1.56$\times$)---solver over-allocation from either indistinguishable distortions (flat) or diluted sensitivity clusters.
    \item \textbf{Mild} (1.07--1.08$\times$): Yi, Llama, Mistral---solvers assign few heads to 2-bit.
\end{itemize}
Notably, Gemma-2 and Qwen-14B reach moderate degradation through different mechanisms (flat over-allocation vs.\ diluted spread compounding), yet the cardinality constraint is effective on both.
This strengthens the ``universal fix'' characterization.

\subsection{Downstream Task Evaluation: Qwen2.5-14B}
\label{sec:downstream_14b}

Table~\ref{tab:14b_downstream} shows 0-shot downstream results on Qwen2.5-14B at 4$\times$ compression.

\begin{table}[t]
\centering
\caption{Qwen2.5-14B downstream task results (0-shot, 4$\times$ compression).}
\label{tab:14b_downstream}
\begin{tabular}{lcccc}
\toprule
Task & BF16 & Uniform 4-bit & MCKP $k{=}12$ & $\Delta$ (pp) \\
\midrule
ARC-Challenge & 0.565 & 0.548 & 0.557 & +0.9 \\
HellaSwag     & 0.800 & 0.785 & 0.793 & +0.8 \\
PIQA          & 0.805 & 0.798 & 0.802 & +0.4 \\
WinoGrande    & 0.715 & 0.705 & 0.712 & +0.7 \\
\midrule
Average       & 0.721 & 0.709 & 0.716 & +0.7 \\
\bottomrule
\end{tabular}
\end{table}

All four tasks show positive deltas (+0.4 to +0.9 percentage points), with 58.3\% average gap recovery.
This is consistent with the PPL improvement: Qwen-7B's 27.3\% PPL gain translates to +2.6pp average downstream, while 14B's 7.9\% PPL gain translates to +0.7pp---roughly proportional.
However, without multi-seed confidence intervals, the per-task deltas (+0.4 to +0.9pp) are within typical single-run noise.
We describe the signal as \emph{modest but consistent across all four tasks}.

\paragraph{Cross-architecture downstream comparison.}
Across prior experiments on five architectures, the diagnostic correctly predicts downstream outcomes:
Qwen-7B (Spread) shows +1.25 to +4.00pp gains across four tasks (72.8\% recovery);
Yi-1.5 (Spike) shows $-0.77$pp average;
Llama (Spike) shows $-0.99$pp average with 0/5 tasks positive---systematic degradation;
Mistral (Flat) shows $-0.48$pp average;
Gemma-2 (Flat) shows $+0.14$pp average (noise range).
ARC-Challenge exhibits the largest per-task sensitivity to allocation strategy across all models (Yi $-3.24$pp, Llama $-2.65$pp), consistent with its reliance on multi-step reasoning where subtle KV cache errors accumulate.

\subsection{Compounding Super-Additivity}
\label{sec:compounding}

We examine the compounding effect by comparing random-$k$ versus top-$k$ 2-bit head selection on Qwen2.5-7B (Table~\ref{tab:compounding}).

\begin{table}[t]
\centering
\caption{Compounding super-additivity: random-$k$ vs.\ top-$k$ 2-bit assignment on Qwen2.5-7B.}
\label{tab:compounding}
\begin{tabular}{rccr}
\toprule
$k$ & top-$k$ PPL & random-$k$ PPL (mean$\pm$std) & Ratio \\
\midrule
1  & 9.826  & 9.887$\pm$0.05  & 1.01 \\
4  & 9.849  & 10.06$\pm$0.06  & 1.02 \\
8  & 10.042 & 11.07$\pm$0.69  & 1.10 \\
16 & 10.42  & 12.91$\pm$0.40  & 1.24 \\
32 & 11.604 & 18.33$\pm$2.53  & 1.58 \\
64 & 22.946 & 55.94$\pm$9.73  & 2.44 \\
\bottomrule
\end{tabular}
\end{table}

Assigning randomly-selected heads to 2-bit produces disproportionately worse perplexity than assigning the highest-sensitivity heads, and this gap widens monotonically with the number of 2-bit heads.
At $k = 64$, random selection yields 2.44$\times$ the degradation of sensitivity-ranked selection.
The variance of the random baseline also grows substantially (std from 0.05 at $k{=}1$ to 9.73 at $k{=}64$), reflecting the increasing influence of which specific heads are forced to extreme quantization.
This super-additive behavior motivates the cardinality constraint: by limiting the number of extreme-quantization assignments, the constraint prevents compounding from low-sensitivity heads that individually appear harmless but collectively degrade quality.
We present this as an empirical observation on a single architecture rather than a formal proposition, and note that cross-architecture characterization of compounding behavior is an open direction.

\subsection{Solver Equivalence Under Constraint}

A consistent finding across all six architectures is that the cardinality constraint eliminates solver dependence.
Table~\ref{tab:solver_equiv} summarizes the unconstrained and constrained solver comparison for the primary showcase (Qwen-7B) and the boundary case (Qwen-14B).

\begin{table}[t]
\centering
\caption{Solver equivalence: greedy vs.\ MCKP, unconstrained vs.\ constrained.}
\label{tab:solver_equiv}
\begin{tabular}{llcc}
\toprule
Model & Setting & Greedy PPL & MCKP PPL \\
\midrule
\multirow{2}{*}{Qwen-7B}  & Unconstrained    & 32.45 & 36.92 \\
                            & Constrained $k{=}8$ & 11.23 & 11.23 \\
\midrule
\multirow{2}{*}{Qwen-14B} & Unconstrained    & 10.15 & 10.35 \\
                            & Constrained $k{=}12$ & 7.18  & 7.18 \\
\bottomrule
\end{tabular}
\end{table}

Without the constraint, solver choice matters (PPL differences of 4.47 and 0.20 respectively).
With the constraint, solver choice is irrelevant ($\Delta = 0.0000$ in both cases).
This pattern holds across all tested architectures, demonstrating that the constraint is the primary intervention---it simultaneously prevents compounding, enables solver interchangeability, and produces the best allocations.

%% ============================================================
%%  4  ANALYSIS
%% ============================================================
\section{Analysis}

\paragraph{Contribution of constraint vs.\ diagnostic.}
Our results support a clear hierarchy: the cardinality constraint is effective on \emph{all} tested architectures (6/6), while the benefit of sensitivity-aware allocation beyond the constraint is architecture-dependent.
On Spread models (Qwen-7B), constrained MCKP achieves 27.3\% PPL improvement.
On diluted-Spread models (Qwen-14B), the benefit is a modest 7.9\%.
On non-Spread models (4/6), uniform 4-bit with the constraint is optimal or near-optimal.
The diagnostic tells practitioners \emph{when} to invest in profiling and non-uniform allocation; the constraint provides \emph{universal insurance} against solver over-allocation.

\paragraph{The role of GQA ratio.}
Yi-1.5-9B (GQA 8:1) has the highest GQA ratio and concentration (0.71) in our study but gains nothing from MCKP ($\approx$0\%).
This definitively refutes the naive hypothesis that higher GQA ratios lead to greater MCKP benefit.
The critical factor is the \emph{shape} of the sensitivity distribution: Yi-1.5's sensitivity is concentrated in a single extreme outlier (L4H3, $\delta = 13.12$, $\tratio = 0.90$---spike pattern) that the solver can handle with a single 8-bit upgrade, leaving no room for multi-head optimization.
Qwen-7B's sensitivity is distributed across a cluster of four heads in Layer~0 ($\tratio = 0.57$---spread pattern) that creates a richer optimization landscape.

\paragraph{Concentration--benefit gradient.}
The six data points trace a non-trivial gradient:
concentration 0.22 $\to$ $-$0.4\% (Mistral),
0.23 $\to$ $-$0.35\% (Gemma-2),
0.47 $\to$ +7.9\% (Qwen-14B),
0.57 $\to$ +0.2\% (Llama),
0.61 $\to$ +27.3\% (Qwen-7B),
0.71 $\to$ $\approx$0\% (Yi-1.5).
The non-monotonicity (Yi-1.5 at 0.71 provides 0\% despite high concentration) demonstrates precisely why $\tratio$ is needed as a second dimension: concentration captures \emph{how much} sensitivity is concentrated, while $\tratio$ captures \emph{how it is concentrated} (cluster vs.\ single outlier).

\paragraph{Universality of the cardinality constraint.}
The constraint prevents degradation through different mechanisms depending on the architecture:
on Spread models, it caps the number of heads entering the compounding regime;
on Flat models (e.g., Gemma-2), it prevents solvers from assigning many heads to 2-bit when per-head distortions are indistinguishable (106 $\to$ 8 heads);
on Spike models, the constraint is non-binding (solvers naturally assign few heads to 2-bit).
In all cases, the effect is beneficial or neutral, never harmful.

%% ============================================================
%%  5  CONCLUSION
%% ============================================================
\section{Conclusion}

We presented a systematic study of per-head KV cache quantization sensitivity across six transformer architectures.
Our central finding is that a cardinality constraint on 2-bit head assignments is a simple, universal intervention that prevents compounding degradation, makes solver choice irrelevant, and yields up to 27.3\% PPL improvement on favorable architectures.

An empirical decision rule based on top-5\% concentration (approximately $\ge 0.50$) and top-1/top-5 ratio (approximately $\le 0.70$) correctly classifies all five non-borderline architectures.
The sixth architecture (Qwen2.5-14B) falls on the decision boundary with a modest 7.9\% benefit, demonstrating that the thresholds are soft boundaries and that benefit degrades continuously with concentration.
We position this rule as a directional heuristic for practitioners, not a definitive diagnostic framework.

\paragraph{Limitations.}
Our study has several limitations.
(1) Only one architecture (Qwen2.5-7B) is a confirmed Spread case with strong MCKP benefit; boundary behavior (14B) provides directional evidence that benefit degrades continuously, but $N = 1$ confirmed Spread limits the generalizability of the strong-benefit claim.
(2) The 14B downstream improvement (+0.7pp average) lacks multi-seed confidence intervals.
(3) The compounding super-additivity observation (Section~\ref{sec:compounding}) is characterized on a single model.
(4) The dilution mechanism linking GQA 5:1 to the diluted spread pattern is a hypothesis supported by one comparison (7B vs.\ 14B within Qwen2.5) but not validated across model families.
(5) We do not evaluate KV cache eviction policies, though our profiling framework extends naturally to that setting.

\paragraph{Future work.}
Cross-architecture compounding characterization, per-layer budget allocation, and integration with KV cache eviction methods are natural extensions.
Validating the decision rule on additional Spread architectures (if they exist) would strengthen the diagnostic from an empirical heuristic to a more robust framework.

\end{document}
